\section{Stability in the common support dictionary}
\label{sec:local-source-stability}
The fixed-scale contraction argument needs convergence of a one-step
map on each fixed input. Its output is a countable list of local
terms. Pointwise convergence of the terms and boundedness of their
total price would not be enough: price could escape to later addresses.
The needed extra property is uniform smallness of the tails. We prove
the summation and local-subtraction argument with this property explicit.
The supports, thresholds and address encodings are fixed throughout.

\subsection{Uniform tails and products}
For nonnegative numbers $a_{N,i}$, indexed by a countable set $I$,
say that their tails are uniformly small if
\[
 \sup_N\sum_{i\in I}a_{N,i}<\infty,
 \qquad
 \forall\epsilon>0\ \exists I_0\subset I\text{ finite}:\quad
 \sup_N\sum_{i\notin I_0}a_{N,i}<\epsilon.
\]
The definition is unchanged if the tail sums are replaced by sums
over arbitrary finite subsets of $I\setminus I_0$. Indeed all
summands are nonnegative. Domination, addition and multiplication
by a fixed nonnegative constant preserve this property. Extension
by zero through a fixed injection preserves it as well.

It also passes to finite-subset products. If
$\sum_i a_{N,i}\le C$, set $A_{N,S}=\prod_{i\in S}a_{N,i}$
for finite subsets $S\subset I$, including the empty subset.
For every finite $J\subset I$,
\[
 \sum_{S\subset J}A_{N,S}
       =\prod_{i\in J}(1+a_{N,i})
       \le\exp\Bigl(\sum_{i\in J}a_{N,i}\Bigr)\le e^C.
\]
This bounds the entire sum by monotone passage through finite $J$.
If $S$ is not contained in $I_0$, choose one of its elements outside
$I_0$ and overcount by all such choices. Removing the chosen element
gives
\[
 \sum_{S\not\subset I_0}A_{N,S}
 \le\sum_{i\notin I_0}a_{N,i}
                   \sum_{T\subset I\setminus\{i\}}A_{N,T}
 \le e^C\sum_{i\notin I_0}a_{N,i}.
\]
Thus the finite collection of subsets of $I_0$ captures all but
an arbitrarily small price. Compatibility restrictions only
remove nonnegative terms and preserve the conclusion. This is
the tail estimate for the finite-subset products in the gas.

There is a corresponding assertion for connected clusters, where
indices may repeat. Suppose the elementary prices $a_{N,i}$ have
uniformly small tails and obey the rooted cluster criterion with
one fixed margin $\rho<1$. The rooted tree bound controls the
total price of clusters of length at least $k+1$ by
\[
                         V\gamma\frac{\rho^k}{1-\rho},
\]
where $V$ is the number of sites. Every cluster has a nonempty
support, so summing the rooted bounds over sites bounds the
unrooted sum. Choose $k$ making this bound small. At each of
the finitely many remaining lengths $r\le k$, the absolute
connected-graph coefficient divided by $r!$ has a finite maximum:
there are finitely many graphs on $r$ labeled vertices. The
remaining price is bounded by a constant times a product of
$r$ elementary prices. The total price of ordered tuples with
some entry outside $I_0$ is at most
\[
                 r C^{r-1}\sum_{i\notin I_0}a_{N,i},
                 \qquad C=\sup_N\sum_i a_{N,i}.
\]
Choose $I_0$ to make these finitely many bounds small. The finite
set of words of length at most $k$ in $I_0$ then captures all
but an arbitrarily small cluster price. This proves uniform
tails for the connected expansion.

For marked clusters the margin is obtained as follows. Suppose
the unmarked criterion has margin $1/2$ and rooted allowance
$\gamma>0$, and the marked elementary prices $v_{N,i}$ have
uniformly small tails and total bound $B$. Let $M$ bound the
support weight in the rooted criterion on this finite torus.
Take
\[
                     \delta=\frac{\gamma}{2(1+MB)}>0.
\]
The additional rooted price of $\delta v$ is at most
$\delta MB\le\gamma/2$. Thus $a+\delta v$ satisfies the
criterion with allowance $3\gamma/2$. The degree price is linear in
this rooted allowance: its original bound is
$\operatorname{degreePrice}(\gamma)\le1/2$, and hence
$\operatorname{degreePrice}(3\gamma/2)
=(3/2)\operatorname{degreePrice}(\gamma)\le3/4$.
For each word, expansion of the product gives
\[
 \sum_{j=1}^r v_{N,i_j}\prod_{h\ne j}a_{N,i_h}
       \le\delta^{-1}\prod_{j=1}^r
                              (a_{N,i_j}+\delta v_{N,i_j}).
\]
The right side includes all terms with exactly one marked factor
and has nonnegative remaining terms. At derivative order $j$,
the product rule costs at most $(r+1)^j\le2^{jr}$; absorb $2^j$
into each elementary price before this comparison. The paid
derivative criterion supplies these enlarged prices. An unmarked
price may be replaced by the smaller of its uniform-tail majorant
and its rooted-criterion majorant, preserving both bounds.
The cluster argument just proved then gives uniform tails for
the differentiated logarithm, not only a bound on its total price.

Let $u_{N,i},v_{N,i}$ take values in a Banach space, and suppose
their norms have uniformly small tails. If
$u_{N,i}-v_{N,i}\to0$ for every fixed $i$, then
\[
                 \sum_i u_{N,i}-\sum_i v_{N,i}\longrightarrow0.
\tag{$\dagger$}\label{eq:source-tight-sum}
\]
To prove it, choose a common finite $I_0$ making each norm tail
less than $\epsilon/4$. The difference of the two sums over
$I_0$ tends to zero, and is eventually less than $\epsilon/2$.
The triangle inequality proves the assertion. Absolute convergence
justifies the subtraction. The same proof applies to supremum norms
of bounded functions and to norms of their fixed-order derivatives.

One useful form of uniform convergence does not require a compact
configuration space. If for every sequence of configurations $U_N$
one has $f_N(U_N)-g_N(U_N)\to0$, then
\[
                    \sup_U|f_N(U)-g_N(U)|\longrightarrow0.
\]
Otherwise some subsequence would have supremum at least $\epsilon$;
choose on it points with difference greater than $\epsilon/2$,
and complete them to a sequence arbitrarily. This contradicts
the hypothesis. The same selection argument applies to derivative
evaluations along normalized local charts. It is this uniform
form, rather than pointwise convergence at each fixed configuration,
that is needed for moving field insertions.

\subsection{From entrywise convergence to a small remainder price}
At the fixed scale let $X_i$ be the regular supports and $Y_i$
the forced supports. Write $M(X)t$ for the support price of an
amplitude $t\ge0$. Since the torus is finite, it has finitely many
supports, and
\[
                         M_*:=\max_X M(X)<\infty.
\]
The price is linear in the amplitude. Suppose two regular lists
$f_{N,i},g_{N,i}$ have derivative majorants with uniformly small
tails, and for every fixed $i$ the difference has arbitrarily
small derivative majorant for sufficiently large $N$.
For any $\epsilon>0$, choose $I_0$ so that the sum of the two
tail majorants is below $\epsilon/(2M_*)$. On $I_0$ use the
common eventual majorant
\[
                     \frac{\epsilon}{2M_*(|I_0|+1)}.
\]
Outside $I_0$ use the sum of the original two majorants.
The resulting majorants bound every difference, and their total
support price is at most $\epsilon$. No supremum over infinitely
many convergence thresholds has been taken.

For forced lists use the identical argument with uniform value
bounds. If their contribution to the remainder norm is divided
by $\kappa>0$, make their total unnormalized support price at
most $\kappa\epsilon/2$, and the regular price at most
$\epsilon/2$. Their sum has remainder seminorm at most
$\epsilon$. This proves smallness in the full remainder space
of Section~\ref{sec:fixed-rung-convergence}. In particular it
does not identify arrays merely because they have the same gas
evaluation.

\subsection{Local subtraction preserves this convergence}
Let $f_{N,i}$ be a regular local term on $X_i$. Its subtraction
uses the fixed linear jet projection $\pi_i$, the fixed local
section $\sigma_i$, and the enlarged support $\widehat X_i$:
\[
 p_{N,i}=\pi_i f_{N,i},\qquad
 r_{N,i}=f_{N,i}-\sigma_i p_{N,i}.
\]
The mask identity leaves the forced rebate
\[
 b_{N,i}
   =\chi_{X_i}f_{N,i}-\sigma_i p_{N,i}
                                  -\chi_{\widehat X_i}r_{N,i}.
\]
At this fixed scale the already established projection and local
section estimates bound each of these three linear maps by a
fixed constant times its input derivative majorant. Subtracting
two inputs gives exactly the same maps of their difference:
\[
 \Delta p_i=\pi_i\Delta f_i,\quad
 \Delta r_i=\Delta f_i-\sigma_i\pi_i\Delta f_i,\quad
 \Delta b_i=\chi_{X_i}\Delta f_i-\sigma_i\pi_i\Delta f_i
                   -\chi_{\widehat X_i}\Delta r_i.
\]
Thus the projected coefficients converge by
\eqref{eq:source-tight-sum}, and the residual and rebate lists
retain the required small-tail bounds and entrywise convergence.
Masks may be discontinuous at the threshold. They are identical
on the two sides and have modulus at most one, so they do not
create a continuity assumption at that threshold.

The forced correction also contains products with the background
activities. On a fixed support--inventory profile it is a sum of
terms
\[
 b_{N,i}(U)\,1_{\mathcal E(i,S,U)}
                            \prod_{j\in S}K_{N,j}(U).
\]
Here $\mathcal E$ is the fixed compatibility and flag condition.
For two backgrounds its indicator is the same: the address
dictionary, supports, inventory and threshold are the same.
For a fixed finite $S$, the identity
\[
 \prod_{j\in S}u_j-\prod_{j\in S}v_j
   =\sum_{j\in S}(u_j-v_j)
          \prod_{k<j}u_k\prod_{k>j}v_k
\]
in any chosen ordering shows convergence of the product when
each activity converges uniformly and is uniformly bounded.
The product-tail estimate above and
\eqref{eq:source-tight-sum} justify summing over $S$. Multiplying
by the convergent rebate, summing the finitely many profiles,
and using the fixed extension of each profile to its address
therefore preserves convergence. The same argument applies to
the raw forced output before the rebate is added.

Consequently local subtraction, barcode grouping, addition of
the background hook, and placement in the common dictionary
preserve convergence of coefficients and smallness of the full
remainder difference. These are exact operations on the common
representation. Changing its forced dictionary midway would
invalidate the product and mask comparisons.

\subsection{Summing the fixed input columns}
For one fixed input remainder, its regular and forced prices
are summable. The selected one-step expansion has one primary
column and one column for each regular and forced input entry.
The primary column has a single fixed bound; the other columns
have bounds proportional to these summable input prices, with
the forced price divided by its fine-scale $\kappa$.
Thus there is a summable sequence $w_i$ bounding, uniformly in
terminal cutoff, both the coefficient norm and remainder price
of output column $i$.

Suppose corresponding individual output columns have the same
regular supports, convergent coefficient differences, and
remainder differences tending to zero in the preceding sense.
Choose a finite set of columns outside which $\sum_i2w_i$
is less than $\epsilon/2$. On that finite set take a common
convergence threshold and sum the small coefficient and remainder
prices. Outside it use $2w_i$. This proves convergence after
all columns have been merged. The unit column is constant and
requires no limiting argument. For the nonunit primary column
the fixed input function is its local-section basis element.
For regular columns it is the fixed masked regular input; for
forced columns it is the fixed profile of the input array.
Hence each column comparison uses the same inserted function
on the two histories, although its background changes.

Signed motions and complex conjugation are bounded linear or
real-linear maps on these data, with the fixed motion prices
already proved. The finite motion average therefore preserves
the convergence. Applying the fixed coefficient pencil and
the common profile capture gives convergence of the actual
one-step coefficient map and its recorded remainder on every
fixed coefficient vector. The coefficient space is finite
dimensional; comparison on its finite basis, followed by the
triangle inequality, gives operator-norm convergence. This is
the local stability input used by the fixed-scale graph argument.

To apply this result to the selected histories one still needs
the individual marked-column comparisons, with uniformly small
tails, for their ordinary and forced expansions. Their hypotheses
are more specific than a uniform total price bound: convergence
of the adjacent path and Gaussian kernels, convergence of the
background activities and coupling, and the common marked input
are all required. The next analytic comparison must supply these
on the same selected histories. The preceding proof isolates
the exact summation, subtraction and normalization steps after
that comparison; it does not replace it by an unproved assertion
of continuity of the infinite expansion.
